% Integration supplement: source recovery, local domain, IFT and compensation.
% Values read from the verified Task 44 JSON outputs; no new calculations here.
\section{Verification, equilibrium framework and compensation}
\label{sec:integrated-context}

\subsection{Production identity and dimensional recovery}
The analysis loads the actual frozen Task 40 objects and WT conserved initial
state. It checks all six stimulated parameter-object hashes and the resting
state hash against the frozen manifest, rather than rebuilding from dataclass
defaults. The source audit compares 150 entries with the Task 40 parent
manifest: 149 match; the sole difference is \texttt{AGENTS.md}, an instructional
metadata change. All listed production Python sources match their recorded
Git blobs. The recovery audit also found no differences in the protected
scientific tree relative to the main head it inspected. The scientific base
is \texttt{c4d4f207f702d8eee09ab94d2d90ae90552dd641}; the recovered Task 44
checkpoint is \texttt{b7b775a0277d93263c4de339745806e5658ec9c2}.
The complete audit and object hashes are retained in
\texttt{output/source\_audit.json}. For an unambiguous initial-state anchor,
the frozen thirteen-state SHA256 is
\begin{center}\footnotesize\ttfamily
a0b96ef16e5967de4dfe2bdabb882f2a28d478e211f920bc19cd735a5a143acb
\end{center}

Dimensional recovery uses $n_{s,i}=C_0V_0x_s$,
$n_{s,l}=C_0L_0y_s$, $V_i=V_0v$, $V_l=L_0\ell$ and $t=t_0\tau$.
If $\widehat I$ is the integrated dimensionless water observable scaled by
$V_0$, then $I=V_0\widehat I$ and $Q=(J_0/C_0)\widehat Q$.
Separate dimensional and dimensionless integrations through 600 s differ by
at most $7.524\times10^{-10}$ in a state component divided by its declared
scale, over the saved comparison times. A separate check recovers cumulative
water to $1.443\times10^{-15}$ pL. The reproduced original sampled WT output
is $0.992524545$ pL, compared with frozen $0.992524544$ pL; its adaptive
integral is $0.992544443$ pL. These two observables have different numerical
definitions. The WT, 5\% expression and null reference checks reproduce
all 34 diagnostic columns at eleven frozen display times, with a largest
cumulative-output discrepancy of $6.74\times10^{-10}$ pL. These are numerical
equivalence checks on the retained equations, not a reduced-model fit.

\subsection{Executable domain and inherited acceptance gates}
The carbonate executable bracket is $[3,11]$ for pH in both compartments.
For the smooth local analysis, require $G(3)<0<G(11)$ for the amount residual
in \eqref{eq:carbonate-amount}, positive transported concentrations and
positive volumes. The NBC electrical solver brackets the basolateral
voltage on $[-0.5,0.5]$ V; strict bracketing of its residual defines the
interior domain used for differentiation. The apical voltage is then
recovered from the coupled current equation, without imposing a separate
apical bracket. The numerical solvers can accept an exact endpoint root;
that fact does not extend the stated smooth interior domain. Require
$V_l>V_{\rm dead}$ to remain on the differentiable, active-outflow branch.
The physical regulatory range is $0\leq r\leq1$; its boundary derivative
uses the inward rule or smooth extension already specified above.

The inherited physiological checks are more restrictive than numerical
speciation. They require all twelve amount/volume components to be strictly
positive and all states, derivatives and diagnostics to be finite, together
with the following intracellular bounds:
\begin{center}
\begin{tabular}{ll}\hline
Quantity & Accepted domain\\\hline
Sodium & $[\mathrm{Na}]_i<40$ mM\\
Potassium & $50\leq[\mathrm K]_i\leq200$ mM\\
Chloride & $30\leq[\mathrm{Cl}]_i\leq80$ mM\\
pH & $6.6\leq\mathrm{pH}_i\leq7.3$\\
Bicarbonate & $0<[\mathrm{HCO}_3]_i<100$ mM\\
Total inorganic carbon & $[T]_i>0$\\
Cell volume & $0<V_i<3$ pL\\\hline
\end{tabular}
\end{center}
The same diagnostics check charge, carbon and buffer accounting to
$10^{-10}$ fmol/s, water accounting to $10^{-12}$ pL/s, each membrane
current residual to $10^{-20}$ A, and compartmental speciation alkalinity
to $10^{-9}$ mM. The declared diagnostic-key set must match, and NBC's
signed charge and carbon sources must equal $-B$ and $2B$ exactly in the
diagnostic representation. These last two sources are not zero residuals.
No additional physiological luminal pH interval has been introduced.

The supplemental full-gate verification independently re-solves 128
parameter-perturbed WT/null roots and 256 nearby expression roots used for
compensation validation. All 384 pass these inherited checks; the largest
conservation residual divided by its own tolerance is $4.086\times10^{-5}$.
The largest dimensionless independent equilibrium residual is
$6.880\times10^{-12}$. These results are retained in
\texttt{output/sensitivity\_full\_gate\_verification.json}; gate acceptance
does not remove the separate thermodynamic qualifications.

\subsection{An eleven-coordinate implicit-function problem}
Let
\[
 z=(x_{\rm Na},x_{\rm K},x_{\rm Cl},x_T,v,
       y_{\rm Na},y_{\rm K},y_{\rm Cl},y_T,\ell,r)\in\mathbb R^{11},
 \qquad \frac{\mathrm dz}{\mathrm d\tau}=F(z,\mu).
\]
The two alkalinity coordinates are reconstructed by
\eqref{eq:neutral-manifold}; both carbon amounts, both volumes and the
regulatory state remain dynamic. The parameter $\mu$ denotes expression
or one specified active model setting, with the scaling fixed at the
reference values. Carbonate and voltage closures are evaluated inside $F$.
Constant full stimulation defines the equilibrium problem $F(z^*,\mu)=0$.
The numerical solver finds ten positive chemical/volume coordinates in
logarithmic variables and sets the \emph{equilibrium} condition $r^*=1$.
All eleven equations are subsequently checked, and all eleven coordinates
are retained in the dynamic Jacobian.

\paragraph{Proposition: local equilibrium sensitivity.}
Suppose the stated interior closure conditions hold, the vector field and
observable $Q(z,\mu)$ are continuously differentiable locally, and
$F_z(z^*,\mu)$ is nonsingular. Then a unique differentiable equilibrium
branch exists in a neighbourhood of this root and parameter. Its derivatives
are
\begin{align}
 \frac{\mathrm dz^*}{\mathrm d\mu}
   &=-F_z^{-1}F_\mu,\label{eq:integrated-ift-state}\\
 \frac{\mathrm dQ^*}{\mathrm d\mu}
   &=Q_\mu-Q_zF_z^{-1}F_\mu.
   \label{eq:integrated-ift-output}
\end{align}
At a physical boundary the claim uses the smooth extension and is restricted
to inward admissible perturbations. The result follows by differentiating
$F(z^*(\mu),\mu)=0$ and applying the implicit function theorem. It establishes
local uniqueness of the branch, not global uniqueness or global existence.
The dimensional dynamic Jacobian is $J=F_z/t_0$; this constant time scaling
does not change \eqref{eq:integrated-ift-state}. $\square$

The terms in \eqref{eq:integrated-ift-output} separate explicit parameter
dependence from the response mediated by the complete state. For AE4
expression the outflow law has $Q_{e_4}=0$. This mathematical zero does not
mean that AE4 has no biological effect on secretion. Likewise, $Q_z$ has
only a lumen-volume component on the active-outflow branch, but every
retained network process can influence that volume.

The expression derivatives were checked against independently solved nearby
equilibria at steps $10^{-4}$ and $5\times10^{-5}$, centred at
$e_4=1,0.5,0.05$ and second-order inward at $e_4=0$. Validation covers all
eleven coordinates and $Q,N,J_4,P,H,\mathrm{pH}_i$. Across these checks,
the largest observable component relative discrepancy is
$1.034\times10^{-7}$; the largest state-derivative relative discrepancy is
$8.070\times10^{-8}$. The maximum nearby-root residual is
$5.683\times10^{-12}$ in the dimensionless independent equations.
This validates local derivatives at the sampled points, without assigning
those difference steps the meaning of uncertainty intervals.

\subsection{Three distinct measures of NKCC compensation}
Because $N$ is an NKCC1 \emph{cycle} flux carrying two chlorides, the local
stationary chloride compensation gain is
\begin{equation}
 C_N(e_4)=-\frac{2\,\mathrm dN^*/\mathrm de_4}
                  {\mathrm dJ_4^*/\mathrm de_4},
 \qquad \mathrm dJ_4^*/\mathrm de_4\ne0.
 \label{eq:integrated-local-compensation}
\end{equation}
Here the parent $J_4$ carries one chloride per cycle. At a zero denominator
the ratio is undefined and the two unnormalised derivatives must be given
separately. The null entry below is an inward derivative and has a nonzero
denominator even though $J_4^*=0$ there.
\begin{center}
\begin{tabular}{rrrr}\hline
$e_4$ & $\mathrm dN^*/\mathrm de_4$ & $\mathrm dJ_4^*/\mathrm de_4$ & $C_N$\\\hline
1 & $-0.0117368$ & $0.0252674$ & 0.929011\\
0.5 & $-0.0283978$ & $0.0620609$ & 0.915159\\
0.05 & $-0.0719960$ & $0.160300$ & 0.898266\\
0 & $-0.0795571$ & $0.177417$ & 0.896838\\\hline
\end{tabular}
\end{center}
Flux derivatives in the table are in fmol/s per unit expression.
Equation \eqref{eq:integrated-local-compensation} describes incremental
replacement at equilibrium; it does not quantify the complete knockout.

For the original finite-time experiment let $\mathcal I=[60,600]$ s and
superscripts $W$ and $0$ denote WT and AE4 null. The relative integrated
NKCC increase and the integrated chloride replacement fraction are instead
\begin{align*}
 R_N&=\frac{\int_{\mathcal I}2(N^0-N^W)\,dt}
             {\int_{\mathcal I}2N^W\,dt}=0.2316345,\\
 C_{N,\mathcal I}&=
 \frac{\int_{\mathcal I}2(N^0-N^W)\,dt}
      {\int_{\mathcal I}(J_4^W-J_4^0)\,dt}=0.8850704.
\end{align*}
These signed integrals use the frozen one-second flux records, with nonzero
denominators. NKCC supplies an additional $36.33701$ fmol of chloride
against $41.05550$ fmol of missing AE4 chloride. Thus $23.16\%$ is an
increase relative to WT NKCC loading, whereas $88.51\%$ is replacement of
lost AE4 chloride over that specified window. Neither is the local gain
$C_N$. The window also differs from the $[0,600]$ s cumulative secretion
comparison. The constant-stimulus stationary null secretion deficit is
$0.94617\%$, compared with the frozen finite-time $3.8575\%$; their
difference retains transient storage and state evolution.

\subsection{Counting local cases and interpreting slow coordinates}
The 26 equilibrium entries are class/expression \emph{cases}: eight C0
expression values and three values for each of the other six source classes.
At zero expression the AE4 source vanishes, so all seven null entries solve
the same system. The count must not be read as 26 distinct equilibria at one
parameter point. Twenty two entries pass the inherited physiological gates.
Across all cases the smallest singular value of the dimensional Jacobian is
at least $9.202\times10^{-4}$ s$^{-1}$ and its condition number ranges from
$1.882\times10^4$ to $2.931\times10^4$ in the declared scaled coordinates.
Nonsingularity supports the local framework; the conditioning and numerical
step checks qualify its implementation.

The slow eigenvectors involve volume and ionic storage together. For the
slowest WT mode (345.82 s), intracellular potassium, cell volume and
intracellular chloride account for $28.82\%$, $20.65\%$ and $18.35\%$ of
the sum of absolute components of the scaled right eigenvector. For the
slowest null mode (459.92 s), the corresponding shares are $31.04\%$,
$24.68\%$ and $16.55\%$. These are coordinate-dependent descriptive shares,
not fractions of physiological energy, transported chloride or causal
control. They supplement the secretion input/output residues and
nonnormality diagnostics reported below. In particular, they give no
justification for removing either volume or declaring the full chemical
system instantaneous.
